\section{Post-training II：偏好优化（DPO/PPO）}

\subsection{从监督信号到偏好信号}
在 SFT 后，模型已具备基本任务执行能力，但答案风格、风险偏好、简洁性与有用性仍需通过偏好学习进一步塑形。讲者围绕 DPO 与 PPO 的实践差异给出了大量工程观察。

\begin{importantbox}{为什么偏好优化对 reasoning 有效}
Reasoning 任务常存在多个可行解，偏好优化能够引导模型朝更可解释、步骤更完整、错误更可纠正的输出分布移动，而不是仅追求单点监督标签匹配。
\end{importantbox}

\begin{figure}[H]
\centering
\includegraphics[width=0.9\textwidth]{frame-05-dpo-ppo.jpg}
\caption{视频约 00:48:30：DPO/PPO 对比及训练稳定性讨论}
\end{figure}

\subsection{DPO 与 PPO 的工程权衡}
\begin{table}[H]
\centering
\renewcommand{\arraystretch}{1.2}
\begin{tabular}{p{2.4cm}p{4.0cm}p{3.8cm}p{3.0cm}}
\toprule
\textbf{方法} & \textbf{优势} & \textbf{风险} & \textbf{适用阶段} \\
\midrule
DPO & 实现相对简单，成本较低，训练链路短 & 偏好样本质量敏感，长度偏置需处理 & 大规模快速迭代 \\
PPO/在线 RL & 可利用在线反馈持续优化，探索更强 & 不稳定、调参复杂、成本高 & 关键能力冲刺阶段 \\
\bottomrule
\end{tabular}
\caption{DPO 与 PPO 的实务差异}
\end{table}

\begin{knowledgebox}{课程里反复出现的一个主题}
算法替换往往不是决定性因素，\textbf{奖励信号质量} 与 \textbf{训练分布覆盖} 通常更关键。很多 ``算法收益'' 实际上来自数据与目标函数修正。
\end{knowledgebox}

\begin{warningbox}{偏好学习的两类隐性退化}
\begin{itemize}
    \item \textbf{过度迎合}：模型学会“看起来像好答案”，但缺少真实推理深度；
    \item \textbf{长度捷径}：模型通过拉长或缩短输出触发 reward 偏好，而非提升正确性。
\end{itemize}
\end{warningbox}

\subsection{本章小结}
偏好优化是后训练中的行为控制层。DPO/PPO 不是二选一，而是要结合数据质量、预算和稳定性目标进行分阶段部署。


